\chapter{Perception}
In this chapter we will look at important facts about human perception to consider when designing diagrams.
An example that teaches this importance can be seen on \url{https://www.csc2.ncsu.edu/faculty/healey/PP/index.html#jscript_search}.

\section{The Human Eye}
Contrary to popular belief, the eye is not a perfect camera.
Instead, it perceives limited colours at different angles \emph{differently}.
Additionally, the brain also has to interpret this information and fill in missing gaps.
A similar conclusion can be made from the \emph{Stroop effect}---where the difference between time needed recognizing meaning of colour and meaning of text delays understanding of colour.
The distribution of cones on the retina makes outer angles appear less detailed then inner angles of the view field.

The so-called \emph{ganglion cells} result in, what is called the \emph{Hermann grid effect}, in which extreme difference in contrast can introduce hallucination.
But in general, these cells enables us perceive contrast and brightness difference between similar colours, when they are directly between each other.
Through them, we notice edges in gradients.
A takeaway from this effect is, that local contrast around edges can be used to make it easier to detect edges.
This is because the human eye perceives \textbf{relative differences} and not absolute differences.

\subsection{Weber's Law \& Fechner's Law}
The quintessence of the last paragraph can be applied through Weber's and Fechner's laws.
Weber's law states that for the notable perceived difference \(K\), stimulus \(S\), and stimulus difference \(dS\) we have \(K \sim \frac{dS}{S}\).
Fechner's Law states that for the perceived brightness \(p\), base stimulus \(S_0\), and sudden current stimulus \(S\), we have \(p \sim \ln \frac{S}{S_0}\).

\subsection{Three Stage Model}
The three stage model divides the processing of visual information into three main phases.

\begin{enumerate}
    \item Rapid pre-attentive processing.
        Goal: Extract low-level properties of a sight.
    \item Pattern perception.
        Goal: Recognizing objects and continuous contours.
    \item Goal-directed processing.
        Driven by attention and linguistic connections.
\end{enumerate}

\noindent
In the first stage, the visual system focuses on the details that \emph{pop out}.
Using pre-attentive features as e.g., marks, exploits this idea.
Pre-attentive features can be: Angle, length, closure, size, curvature, density, count, colour hue, colour intensity, intersection, termination, 3d depth cue, lighting direction, flicker, motion direction, motion velocity.

It has to be noted, that the overuse or random use of pre-attentive features can minder this effect, and make it even more complicated to highlight e.g., points.

\begin{center}
    \todo[inline]{TODO: Maybe fill in missing content}
\end{center}

\section{Gestalt Theory}
In Gestalt theory, the law of \emph{Prägnanz} derives different laws based on how the mind groups shapes together to form objects.
There are different laws based on this principle\footnote{To learn and differentiate between each law, I would recommend looking at the pictures in the slides and find external examples, e.g., \url{https://en.wikipedia.org/wiki/Principles_of_grouping}.}.

\begin{description}
    \item[Good Law of Prägnanz] Ambiguous and complex images are perceived simplified.
    \item[Law of Proximity] Visually close things are related.
    \item[Law of Similarity] Things that look alike are related.
    \item[Law of Connection] Visually connected things are related.
    \item[Law of Closure] Incomplete shapes are interpreted as complete.
    \item[Law of Enclosure] Enclosed shapes are related.
    \item[Law of Continuity] Hidden objects are interpreted as fitting simple shapes.
    \item[Law of Symmetry] Symmetric things are related.
    \item[Law of Past Experience] Shapes that form known pictures are related.
    \item[Law of Figure and Ground] Elements are either perceived as figures or background.
    \item[Law of Common Fate] Elements moving in the same direction are related.
\end{description}

